\documentclass[12pt]{article}
\usepackage[utf8]{inputenc}
\usepackage[T1]{fontenc}
\usepackage{hyperref}
\title{Yela Language Resources for K-12 Curriculum Development}
\author{Compiled by Claude Sonnet 5.5 (claude-sonnet-5-5) for Andr\'as Kornai}
\date{}
\begin{document}
\maketitle

\section{Introduction}
Yela (ISO 639-3 \texttt{yel}, Glottocode yela1238) is a code for a language in the Democratic Republic of the Congo. The task description states that the Ethnologue editors found no evidence that it exists as a distinct language. My searches found no confirmation of a retirement. They did find a related picture: Wikipedia's Yela-Kela article treats the name as a Bantu language of {\'E}quateur Province with two codes, \texttt{kel} (Kela) and \texttt{yel} (Yela). Glottolog notes a linguist's view that the two are not clearly distinct and might be merged. The status of \texttt{yel} remains unclear, and I could not settle it.

\section{Primary Educational Resources}
None were found.

\section{Linguistic and Reference Materials}
\begin{itemize}
\item Wikipedia article on Yela-Kela: \url{https://en.wikipedia.org/wiki/Yela-Kela_language}
\item Glottolog entry for ISO code \texttt{yel}: \url{https://glottolog.org/resource/languoid/iso/yel}
\end{itemize}

\section{Religious and Community Resources}
None were found.

\section{Digital and Audio Resources}
None were found.

\section{Curriculum Development Context}
No curriculum resources exist for Yela. If the code is a duplicate of Kela, any work should be planned for Kela and its neighbours in {\'E}quateur Province. That is a question to settle with the registry and the community before any materials are produced.
\end{document}